\subsection{Specialization of the corrected tangent lattices}
\label{sec:geometry-hodge-specialization}

The correction used in the point-height argument is made on the original
normal compactification, before specializing to a boundary point. This
order matters. A semiabelian variety of dimension $g$ and torus rank $r$
has $p$-adic Tate module of rank $2g-r$. Its intrinsic torsion cannot be
substituted for the rank-$2g$ torsion of the generic abelian family.
We prove the required comparison by applying the finite-group
Hodge--Tate theorem on the original charts and then specializing the
original locally free modules.

Let $X$ be one of the original normal integral models, with function
field $K$ of characteristic zero, and let $G/X$ be its smooth
semiabelian family of positive relative dimension. Its generic fibre
is abelian. Retain its coefficient
action by $\mathcal O_E$, its rigidified cubical line bundle, and its
generic principal polarization. Write
\[
 \mathcal T=\operatorname{Lie}(G/X),\qquad
 \omega=\mathcal T^\vee.
\]
All coefficient embeddings used below are the original labels. Their
values on $\mathcal O_E$ are regular: they are integral over $\mathbb Z$
in the function field of a normal model.

For $p\in S$, put $n_p=p^{R_p}$, where the prescribed $R_p\ge4$,
and let $T_p=G_K[n_p](K)$. The original level makes all geometric
$n_p$-torsion rational over $K$. For $t\in T_p$, differentiate at the
identity the original Weil character $e_{n_p}(t,-)$ on $G[n_p]$.
As proved below, this gives a canonical local section
\[
 \delta_t\in\omega/n_p\omega.
\]
Define $\mathcal N_p\subset\omega$ to be the inverse image of the
span of these sections. Thus it contains $n_p\omega$, and on a chart
it is generated by arbitrary lifts of the $\delta_t$, together with
$n_p\omega$. This is precisely the cotangent module used in the
original construction; changing a lift changes it by $n_p\omega$.

Let $\mathcal T_{K,\sigma}$ be the $\sigma$-eigenspace of the
coefficient action on $\mathcal T_K$. Set
\[
 \mathcal Q_{p,\sigma}
   =\operatorname{im}\bigl(\mathcal N_p\longrightarrow
                          \mathcal T_{K,\sigma}^{\vee}\bigr),
 \qquad
 \mathcal D_{p,\sigma}=\mathcal Q_{p,\sigma}^{\vee}
       \subset\mathcal T_{K,\sigma}.
 \tag{HS1}
\]
We assume the original local-freeness condition: $\mathcal N_p$ and
each $\mathcal Q_{p,\sigma}$ are locally free of their prescribed
ranks near the fibre over $p$. In the selected half of the labels the
corrected module is
$\sigma(\mathfrak D_E)^{-1}\mathcal D_{p,\sigma}$; in the other
half it is $\mathcal D_{p,\sigma}$. Here $\mathfrak D_E$ is the
different of $E/\mathbb Q$, extended to the chart. Away from $S$ the
original corrected bundle is the ordinary integral eigenbundle.
Denote this globally prescribed bundle by $\mathcal F_\sigma$.

\begin{theorem}[The actual corrected fibre]
\label{thm:actual-hodge-specialization}
Suppose $S$ contains every rational prime dividing
$\operatorname{disc}(E)$. Let $z:\operatorname{Spec}\mathcal O_L\to X$
be an original arithmetic point, including a boundary point, with
its original coefficient labels $\sigma:E\hookrightarrow L$.
Let
\[
 \Lambda_\sigma
   =\operatorname{Lie}(G_z/\mathcal O_L)
                   \cap\operatorname{Lie}(G_{z,L})_\sigma,
 \qquad
 F_\sigma=(z^*\mathcal F_\sigma)
                   \subset\operatorname{Lie}(G_{z,L})_\sigma.
\]
The inclusion defining $F_\sigma$ is the original evaluation of
sections of the corrected bundle. There is a nonzero integral ideal
\[
 \mathfrak a_\sigma
   =\sigma(\mathfrak D_E)\mathcal O_L
                   \prod_{p\in S}(\kappa_{p,z}),
 \qquad
 \kappa_{p,z}=
 \begin{cases}
  4,&p=2,\\
  \zeta_{p,z}-1,&p\ne2,
 \end{cases}
 \tag{HS2}
\]
where $\zeta_{p,z}$ is the specialization of an original primitive
$p$th root of unity, such that
\begin{equation}\label{eq:actual-hodge-inclusions}
 \mathfrak a_\sigma F_\sigma\subset\Lambda_\sigma\subset F_\sigma.
\end{equation}
Consequently $\mathfrak a_\sigma\Lambda_\sigma\subset F_\sigma$ as
well. Its exact normalized norm is
\begin{equation}\label{eq:actual-hodge-ideal-norm}
 \frac{\log N_{L/\mathbb Q}\mathfrak a_\sigma}{[L:\mathbb Q]}
 =\log\operatorname{rd}(E)
  +\sum_{\substack{p\in S\\p\ne2}}\frac{\log p}{p-1}
  +2\,\mathbf 1_{2\in S}\log2
 \le 2\log\operatorname{rd}(E)
       +2\sum_{p\in S}\frac{\log p}{p-1}.
\end{equation}
Every lattice and label here is the original one. In particular,
the assertion applies without replacing an image lattice in an
abelian quotient by its saturation.
\end{theorem}

\begin{proof}
We separate the finite-group estimate from the specialization argument.

\emph{The original characters and their differentials.}
Work first over a normal affine chart $A\subset K$ of $X$ and put
$n=n_p$. The kernel $\Gamma=G[n]$ is separated, flat and quasi-finite
of finite presentation over $A$. At a point of the chart pass to
the henselization of its normal local ring. The finite-part theorem
for a separated quasi-finite scheme over a henselian local base
gives an open and closed finite subscheme $\mathcal H\subset\Gamma$
containing the whole closed fibre\cite{hodge-finite-part}. It is a
subgroup: multiplication and inversion preserve it on the closed
fibre, and the inverse image of the complementary open and closed
part in a finite scheme is empty if its closed fibre is empty.
It is finite flat because $\Gamma$ is flat.

The restriction of $e_n(t,-)$ to $\mathcal H_K$ is a rational point
of the finite Cartier dual $\mathcal H^D$. It extends uniquely over
the normal henselian ring: the images of its finite coordinate
algebra are integral over the base and lie in its fraction field.
The resulting character is regular near the identity of $\Gamma$.
Its generic rational function was defined before henselization;
faithfully flat descent of membership in the local ring therefore
gives a regular germ on the original $\Gamma$ at the identity.
Uniqueness follows from flatness and generic-fibre density.

Lift this germ through the closed immersion $\Gamma\hookrightarrow G$.
The cotangent sequence at the identity identifies
\[
 \omega_\Gamma=\operatorname{coker}
       \bigl([n]^*:\omega\longrightarrow\omega\bigr)
       =\omega/n\omega,
 \tag{HS3}
\]
since the differential of multiplication by $n$ is multiplication
by $n$. The differential of the lift is consequently independent
of its choice modulo $n\omega$. Its value is $\delta_t$.
The character has value $1$ at the identity, so its ordinary
differential there is its logarithmic differential.

This also checks the definition of $\mathcal N_p$ against the
original one: a regular function realizing the original generic
Weil character has differential congruent to $\delta_t$ modulo
$n\omega$, and every lift of $\delta_t$ occurs in this way at the
level of differentials. The zero character supplies $n\omega$.
Thus no new characters or tangent module have been introduced.
For an injective valued extension of $A$, the characters agree
after base change by their generic values and flatness. Their
differentials and the generated module agree as well. We use this
compatibility only for injective test extensions; specialization
with a kernel will be handled separately below.

\emph{A finite flat model of the full generic torsion.}
Let $C$ be a complete algebraically closed rank-one valued extension
of $K$, with valuation ring $\mathcal O_C$ and $v(p)=1$. The
generic fibre of $G_{\mathcal O_C}$ is still abelian. Put
$T=G_C[n](C)$ and let $\mathcal H\subset G_{\mathcal O_C}[n]$ be
its finite part containing the closed fibre. Write $H=\mathcal H(C)$.
Properness identifies $H$ with $\mathcal H(\mathcal O_C)$.

There is a finite flat commutative group scheme $\mathcal B$ whose
generic group is all of $T$: take one copy of $\mathcal H$ for each
coset of $H$ in $T$. Choose representatives $t_c$ with $t_0=0$ and define the sum
on these copies by
\[
 (c,h)+(d,h')=(c+d,h+h'+t_c+t_d-t_{c+d}).
 \tag{HS4}
\]
The last term belongs to $H$ and hence is an integral section of
$\mathcal H$. The cocycle identity proves associativity. Inversion
and the identity are defined by the same formula. The resulting
scheme is independent of representatives by translation, is killed
by $n$, and is finite flat of rank $|T|=n^{2\dim G_K}$.
The identity copy is $\mathcal H$, so
\[
 \omega_{\mathcal B}=\omega_{\mathcal H}
                      =\omega_{G_{\mathcal O_C}}/n\omega_{G_{\mathcal O_C}}.
 \tag{HS5}
\]
This is also the original induced finite Hopf order: a function on
$T$ belongs to it exactly when its restriction to every translate
of $H$ is regular on $\mathcal H$.

Here the original cubical line supplies a necessary integral
compatibility. If $\mathcal L$ is that rigidified line, set
\[
 \mathcal P(x,y)=\mathcal L_{x+y}\otimes\mathcal L_x^{-1}
                     \otimes\mathcal L_y^{-1}\otimes\mathcal L_0.
\]
The cubical structure gives its biextension laws: addition in
either variable tensors the corresponding lines, associatively
and compatibly with addition in the other variable. The
rigidification trivializes $\mathcal P(0,y)$ and
$\mathcal P(x,0)$. On $G[n]\times G[n]$, the two $n$-fold
addition laws consequently give two nowhere-vanishing
trivializations of the same line $\mathcal P(x,y)^{\otimes n}$,
one through $\mathcal P(nx,y)$ and one through
$\mathcal P(x,ny)$. Their ratio, in the order of the original
Weil convention, is a regular unit. This construction is made
with the universal torsion points and commutes with every base
change; it is a regular function on the product scheme, not only
a collection of functions at its rational points.

To check the normalization, choose lifts of $x$ and $y$ to the
theta group of $\mathcal L^{\otimes n}$ on a geometric generic
fibre. The cubical addition law writes the two translation
isomorphisms as the corresponding $n$-fold trivializations,
each multiplied by a scalar. In their ordered commutator the two
scalars cancel, leaving exactly their ratio. Hence the resulting
unit is the original $e_n(x,y)$, rather than an unspecified power
or scalar multiple of it. Its bilinear identities hold on the
generic fibre and extend by flatness. Restricting this construction
gives the required regular pairing on
$\mathcal H\times\mathcal H$\cite{hodge-specialization-source}.
For arbitrary $t,t'\in T$,
the mixed characters $e_n(t,-)|_H$ and $e_n(-,t')|_H$ extend over
$\mathcal H$, since their points in $\mathcal H^D(C)$ extend by
properness. On two coset copies the formula
\[
 e_n(t+h,t'+h')
   =e_n(t,t')e_n(t,h')e_n(h,t')e_n(h,h')
 \tag{HS6}
\]
therefore defines a regular pairing. The constant factor is a
root of unity and hence an integral unit. The identities needed
to glue and to make the pairing bilinear hold generically and
thus hold on the finite flat schemes. We have extended exactly
the original pairing to a homomorphism
$\lambda:\mathcal B\to\mathcal B^D$.

The original generic principal polarization makes $\lambda_C$
an isomorphism. We do not need to claim that $\lambda$ is an
integral isomorphism. The maps from integral points to $C$-points
of both finite proper schemes are bijections. It follows that
\[
 \lambda:\mathcal B(\mathcal O_C)
                       \longrightarrow\mathcal B^D(\mathcal O_C)
 \quad\hbox{is bijective}.
 \tag{HS7}
\]

\emph{The finite-group Hodge--Tate bound.}
For any commutative finite flat group $J/\mathcal O_C$ of $p$-power
order, Fargues's finite-group theorem says that the cokernel of
\[
 J(\mathcal O_C)\otimes_{\mathbb Z}\mathcal O_C
       \longrightarrow\omega_{J^D}
\]
is annihilated by every element of valuation at least $1/(p-1)$
\cite[\S5.4, Theorem~3]{hodge-fargues-finite}. His base-field
convention allows complete rank-one valuations that are not discrete
and imposes no condition on the residue field. The theorem includes
$p=2$; we retain the original weaker annihilator $4$ in that case.
Apply it to $J=\mathcal B^D$. By \textup{(HS5)--(HS7)}, its
Hodge--Tate differentials are exactly the original $\delta_t$ in
$\omega/n\omega$. With
$\kappa_p=\zeta_p-1$ for odd $p$ and $\kappa_2=4$, it gives
\[
 \kappa_p\omega_{\mathcal O_C}
             \subset\mathcal N_{p,\mathcal O_C}
             \subset\omega_{\mathcal O_C}.
 \tag{HS8}
\]
Here $v(\zeta_p-1)=1/(p-1)$ and $v(4)=2$ at $p=2$.

Return to the normal noetherian chart $A$. The generic full level
and the perfect generic pairing supply the chosen primitive roots,
which are integral on $A$ by normality. At a height-one prime over
$p$, embed the fraction field of its completed discrete valuation
ring into a complete algebraically closed valued field. This is an
injective test extension of $A$, so the preceding identification
of $\mathcal N_p$ applies. The extension of valuation rings is
faithfully flat. Thus \textup{(HS8)} descends to $A_{\mathfrak q}$.
At height-one primes not over $p$ we already have
$\mathcal N_p=\omega$, since $n$ is a unit. Near the fibre over
$p$, local freeness makes $\mathcal N_p$ reflexive. Intersecting
at height-one primes in a normal ring proves the original chart
inclusions
\begin{equation}\label{eq:hodge-chart-inclusions}
 \kappa_p\omega\subset\mathcal N_p\subset\omega.
\end{equation}
This argument has used a single finite level. It requires neither
a compatible tower of finite flat groups nor a replacement of the
given boundary point by an abelian variety.

\emph{Specializing the restricted dual.}
We give the algebra needed when $z$ has a kernel on a chart.
Let $A\to R$ be the resulting map to a characteristic-zero
valuation ring, let $C_R=\operatorname{Frac}(R)$, and write
$M=\mathcal T(A)$ and $N=\mathcal N_p(A)$ after shrinking so the
relevant modules are free. For the fixed label put
\[
 Q=\operatorname{im}(N\longrightarrow M_{K,\sigma}^{\vee}),
 \qquad D=Q^\vee.
\]
The original local-freeness condition makes $Q$ projective. Hence
$N\twoheadrightarrow Q$ is a split surjection and
$D\otimes_A R=(Q\otimes_A R)^\vee$.
The chart inclusions imply, in the original generic eigenspace,
\[
 M\cap M_{K,\sigma}\subset D,\qquad \kappa_p D\subset M.
 \tag{HS9}
\]

To identify the specialized subspace, choose a nonzero integer
$b_E$ clearing the denominators in the trace-dual integral basis
of $E$. The $\sigma$-projector, expressed using that dual basis,
is a polynomial in the coefficient action with coefficients in
$A[1/b_E]$; the same formula specializes to the $\sigma$-projector
over $C_R$. After inverting $b_E\kappa_p$, \textup{(HS9)}
identifies $D$ with that eigensummand of $M$. Both constants
specialize to nonzero elements of $C_R$. Thus the map obtained by
specializing $\kappa_p D\to M$ and dividing by $\kappa_{p,z}$
identifies $D\otimes_A C_R$ with the actual specialized
$\sigma$-eigenspace. Since $D\otimes_A R$ is projective, it
embeds in that space. This is the evaluation of the original
sections, because it agrees after multiplication by the same
nonzero scalar with evaluation of integral tangent sections.

Put $\Lambda_{\sigma,R}=M_R\cap M_{C_R,\sigma}$. For
$v\in\Lambda_{\sigma,R}$, evaluation on $N_R$ is integral. It
annihilates $\ker(N\to Q)\otimes_A R$: after passing to $C_R$
this kernel is in the annihilator of the same eigenspace, by the
projector formula, and an element of $R$ vanishing in $C_R$ is
zero. Evaluation therefore descends to an element of
$(Q\otimes_A R)^\vee=D\otimes_A R$. Under the preceding rational
identification this element is $v$ itself. Conversely
$\kappa_{p,z}(D\otimes_A R)$ lies in $M_R$ by \textup{(HS9)}
and remains in the same eigenspace. We obtain
\begin{equation}\label{eq:hodge-dual-specialization}
 \kappa_{p,z}D_R\subset\Lambda_{\sigma,R}\subset D_R.
\end{equation}
The integer $b_E$ was used only to verify identities over $C_R$.
It contributes no integral multiplier to this inclusion.

The embedded different is an invertible integral ideal, so its
inverse commutes with this specialization. For either original
half of the labels, the definition of $\mathcal F_\sigma$ and
\eqref{eq:hodge-dual-specialization} give
\[
 \kappa_{p,z}\sigma(\mathfrak D_E)F_{\sigma,R}
                 \subset\Lambda_{\sigma,R}\subset F_{\sigma,R}.
 \tag{HS10}
\]
At a finite place outside $S$, the discriminant is a unit. The
coefficient algebra is then split etale on the coefficient cover,
and its projectors are integral. The ordinary eigenbundle is a
direct summand and commutes with specialization. Hence
$F_{\sigma,R}=\Lambda_{\sigma,R}$ there.
These are fibres of the original locally free bundles, so their
local section evaluations are exactly the localizations of the
global lattices in the theorem. Intersecting at all finite places
of $L$ proves \eqref{eq:actual-hodge-inclusions} with the ideal
\textup{(HS2)}.

\emph{The ideal norm.}
The specialized odd root is still primitive: its value satisfies
the cyclotomic polynomial $\Phi_p$, and the target has
characteristic zero. Thus every $\kappa_{p,z}$ is nonzero. The
cyclotomic identity gives
\[
 \frac{\log N_{L/\mathbb Q}(\zeta_{p,z}-1)}{[L:\mathbb Q]}
       =\frac{\log p}{p-1},\qquad
 \frac{\log N_{L/\mathbb Q}(4)}{[L:\mathbb Q]}=2\log2.
\]
Extension of the different gives
\[
 \frac{\log N_{L/\mathbb Q}
                  (\sigma(\mathfrak D_E)\mathcal O_L)}{[L:\mathbb Q]}
       =\frac{\log|\operatorname{disc}(E)|}{[E:\mathbb Q]}
       =\log\operatorname{rd}(E).
\]
Multiplicativity proves the equality in
\eqref{eq:actual-hodge-ideal-norm}; its upper bound follows from
$\operatorname{rd}(E)\ge1$ and nonnegativity of the prime terms.
\end{proof}

\paragraph{Passage to the abelian quotient.}
Let $q:G_{z,L}\to A_z$ be the original abelian quotient, and retain
its differential on the original lattices. Distinguish
\[
 W_\sigma=\text{the full integral eigenlattice of }A_z,\qquad
 V_\sigma=q_*\Lambda_\sigma,\qquad
 \overline F_\sigma=q_*F_\sigma.
\]
The theorem gives
\[
 \mathfrak a_\sigma\overline F_\sigma
                       \subset V_\sigma\subset\overline F_\sigma.
 \tag{HS11}
\]
It gives no identification of $V_\sigma$ with $W_\sigma$.
The finite index $[W_\sigma:V_\sigma]$ in the actual point-height
bound remains unchanged. If $r_\sigma=\dim_L\operatorname{Lie}(A_z)_\sigma$
and both image lattices carry the prescribed common quotient
metric, then
\[
 0\le\widehat{\deg}_{\rm n}\overline F_\sigma
          -\widehat{\deg}_{\rm n}V_\sigma
 \le r_\sigma\,
      \frac{\log N_{L/\mathbb Q}\mathfrak a_\sigma}{[L:\mathbb Q]}.
 \tag{HS12}
\]
Indeed the quotient $\overline F_\sigma/V_\sigma$ is annihilated
by $\mathfrak a_\sigma$ and is generated locally by at most
$r_\sigma$ elements. This is the finite correction needed by the
point-height comparison. The regularized intersection identity
is a separate obligation.

\begin{proposition}[Simultaneous local freeness on the original model]
\label{prop:actual-hodge-free-model}
Start with the prescribed normal projective model of finite type
over $\mathbb Z$, the original full levels $p^{R_p}$ for the finite
set $S$, the same generic coefficient field, and the same
semiabelian family and cubical line. There is a normal projective
birational modification $f:X'\to X$, isomorphic over the whole
characteristic-zero fibre, on which all the actual modules
$\mathcal N_p$ and $\mathcal Q_{p,\sigma}$ are locally free.
The original labels, levels, generic field, and generic family
are unchanged. The construction gives precisely the required
corrected bundles on $X'$ and hence supplies the local-freeness
hypothesis of Theorem~\ref{thm:actual-hodge-specialization}.

If some of these modules are already locally free, they are
pulled back as the same locally free modules throughout the
construction. Every original arithmetic point lifts uniquely
to $X'$, because $f$ is an isomorphism on its characteristic-zero
generic point. The comparison in the theorem concerns the actual
modules at this lift and the pullback of the same semiabelian
family.
\end{proposition}

\begin{proof}
We first verify the compatibility that an abstract flattening
argument would not supply. For a dominant morphism $f:X'\to X$
between normal integral models with the same generic field,
the pullback of a regular germ realizing $e_{n_p}(t,-)$ is a
regular germ realizing that same character on the pulled-back
family. Any other such germ on its torsion kernel has the same
generic restriction, and equality follows from flatness over
the integral base. Their differentials modulo $n_p\omega$ are
therefore identical. Since the set of original generic torsion
labels is unchanged, and the original smooth cotangent bundle
commutes with base change, we obtain the equality of actual
modules
\[
 \mathcal N'_p
   =\operatorname{im}(f^*\mathcal N_p\longrightarrow f^*\omega).
 \tag{HS13}
\]
Restriction to the same rational eigenspace gives, with the
original coefficient projector,
\[
 \mathcal Q'_{p,\sigma}
   =\operatorname{im}(f^*\mathcal Q_{p,\sigma}
                      \longrightarrow\mathcal T_{K,\sigma}^{\vee}).
 \tag{HS14}
\]
Thus these modules are their torsion-free pullbacks, not arbitrary
coherent images on the modified model. The displayed equalities
are the original Weil-differential pullback identities.

The local description by finitely many $\delta_t$ and
$n_p\omega$ proves coherence of $\mathcal N_p$; each
$\mathcal Q_{p,\sigma}$ is coherent as the image of a coherent
module in a fixed rational vector space. After inverting $p$,
$\mathcal N_p=\omega$. Over the characteristic-zero fibre, the
coefficient action splits into its eigensummands, so each
$\mathcal Q_{p,\sigma}$ is the corresponding locally free
cotangent summand. Their ranks are the prescribed signature
ranks. Hence all these finitely many coherent modules have
constant generic rank and are locally free on the entire
characteristic-zero fibre.

For a coherent module of generic rank $a$, blow up its $a$th
Fitting ideal. Its strict transform is locally free of rank $a$
\cite[Lemma~31.36.3]{hodge-flattening}. On an integral dominating
model the strict transform is the image in the original rational
vector space: the kernel of the map from the pulled-back module
is exactly its torsion, supported over the exceptional locus.
Normalize this blowup. Normalization is finite because the models
are excellent, and it preserves the locally free strict transform.
Equations \textup{(HS13)--(HS14)} identify this transform with
the actual module calculated from the pulled-back semiabelian
family and its original Weil characters.

Repeat for the finite list of modules. A module already made
locally free stays so: its pullback is locally free and the map
to the original rational vector space is injective, since it is
injective generically on an integral model. Every centre is
vertical, because all the modules are locally free in
characteristic zero. The composite is therefore projective,
birational, normal after the indicated normalizations, and an
isomorphism over the whole characteristic-zero fibre. It remains
flat over $\mathbb Z$: its integral affine coordinate rings have
characteristic zero and are torsion-free over $\mathbb Z$.
Rank-zero
restricted modules are zero and require no modification.

The same generic field and torsion labels remain throughout;
pullback retains smoothness, the semiabelian structure, the
coefficient action, and the rigidified cubical line and its
principal generic polarization. The different is an invertible
coefficient ideal, so dualizing the locally free
$\mathcal Q'_{p,\sigma}$ and applying exactly the original
different correction gives locally free corrected modules.
On the open subset where $p$ is invertible and the other primes
of $S$ have been removed, these agree with the original ordinary
eigenbundle: its coefficient projectors are integral there.
The open subsets obtained by removing all primes of $S$ except
one, together with the complement of $S$, cover $X'$. They glue
these actual modules to the required $\mathcal F'_{\sigma}$.

The construction retains the old projective coordinates by
pullback; adjoining homogeneous coordinates for the projective
modification realizes it as the permitted refinement of the
original model. Every previously prescribed line bundle and
metric is pulled back, with all its original finite factors.
The complex fibre and its metrics are unchanged. We make no
additional numerical claim about the added projective
coordinates. Finally, the valuative criterion for properness
lifts an original arithmetic point at each finite place, from
its unique lift in characteristic zero. Separatedness makes the
lifts agree and glue over $\mathcal O_L$. No extension of $L$
is required.
\end{proof}
